% ============================================================
%  CHAPITRE 4 : RÉALISATION
% ============================================================
\chapter{Réalisation}

\section*{Introduction}
\addcontentsline{toc}{section}{Introduction}
Ce chapitre présente la mise en œuvre technique de la plateforme web interactive dédiée à la valorisation de l'huile d'olive FENDRI. Il détaille les différentes interfaces développées, les fonctionnalités implémentées ainsi que les choix techniques adoptés pour garantir une expérience utilisateur fluide et immersive.

% -----------------------------------------------------------
\section{Page d'accueil}
% -----------------------------------------------------------

\subsection{Section héro}
% [Copier depuis Canva]
\begin{figure}[H]
  \centering
  \includegraphics[width=\textwidth]{hero.png}
  \caption{Section héro de la page d'accueil}
  \label{fig:hero}
\end{figure}

\subsection{Section Nos Huiles}
% [Copier depuis Canva]
\begin{figure}[H]
  \centering
  \includegraphics[width=\textwidth]{section_huiles.png}
  \caption{Section Nos Huiles}
  \label{fig:section_huiles}
\end{figure}

\subsection{Section Notre Usine}
% [Copier depuis Canva]
\begin{figure}[H]
  \centering
  \includegraphics[width=\textwidth]{section_usine.png}
  \caption{Section Usine}
  \label{fig:section_usine}
\end{figure}

\subsection{Section Engagements}
% [Copier depuis Canva]
\begin{figure}[H]
  \centering
  \includegraphics[width=\textwidth]{section_engagements.png}
  \caption{Section Engagements}
  \label{fig:section_engagements}
\end{figure}

\subsection{Section Récompenses}
% [Copier depuis Canva]
\begin{figure}[H]
  \centering
  \includegraphics[width=\textwidth]{section_recompenses.png}
  \caption{Section Récompenses}
  \label{fig:section_recompenses}
\end{figure}

\subsection{Section Contact}
% [Copier depuis Canva]
\begin{figure}[H]
  \centering
  \includegraphics[width=\textwidth]{section_contact.png}
  \caption{Section Contact}
  \label{fig:section_contact}
\end{figure}

\begin{figure}[H]
  \centering
  \includegraphics[width=\textwidth]{confirmation_envoi.png}
  \caption{Confirmation d'envoi}
  \label{fig:confirmation_envoi}
\end{figure}

\subsection{Section FAQ}
% [Copier depuis Canva]
\begin{figure}[H]
  \centering
  \includegraphics[width=\textwidth]{faq.png}
  \caption{Page FAQ}
  \label{fig:faq}
\end{figure}

% -----------------------------------------------------------
\section{Page collection}
% -----------------------------------------------------------
% [Copier depuis Canva]
\begin{figure}[H]
  \centering
  
\includegraphics[width=\textwidth]{collection.png}
  \caption{Page collection}
  \label{fig:collection}
\end{figure}

\begin{figure}[H]
  \centering
  \includegraphics[width=\textwidth]{fiche_produit.png}
  \caption{Fiche de produit}
  \label{fig:fiche_produit}
\end{figure}

\subsection{Panier et commande}
% [Copier depuis Canva]
\begin{figure}[H]
  \centering
  \includegraphics[width=\textwidth]{panier.png}
  \caption{Panier}
  \label{fig:panier}
\end{figure}

\begin{figure}[H]
  \centering
  \includegraphics[width=\textwidth]{validation_commande.png}
  \caption{Remplissage et validation des champs}
  \label{fig:validation_commande}
\end{figure}

% -----------------------------------------------------------
\section{Configurateur}
% -----------------------------------------------------------
% [Copier depuis Canva]

\subsection{Étape 1 — Choix du modèle de bouteille}
\begin{figure}[H]
  \centering
  \includegraphics[width=\textwidth]{config_etape1.png}
  \caption{Étape 1 du configurateur}
  \label{fig:config_etape1}
\end{figure}

\subsection{Étape 2 — Sélection de la contenance}
\begin{figure}[H]
  \centering
  \includegraphics[width=\textwidth]{config_etape2.png}
  \caption{Étape 2 du configurateur}
  \label{fig:config_etape2}
\end{figure}

\subsection{Étape 3 — Sélection de l'étiquette}
\begin{figure}[H]
  \centering
  \includegraphics[width=\textwidth]{config_etape3.png}
  \caption{Étape 3 du configurateur}
  \label{fig:config_etape3}
\end{figure}

\subsection{Étape 4 — Sélection de l'emballage}
\begin{figure}[H]
  \centering
  \includegraphics[width=\textwidth]{config_etape4.png}
  \caption{Étape 4 du configurateur}
  \label{fig:config_etape4}
\end{figure}

\subsection{Étape 5 — Personnalisation d'un message}
\begin{figure}[H]
  \centering
  \includegraphics[width=\textwidth]{config_etape5.png}
  \caption{Étape 5 du configurateur}
  \label{fig:config_etape5}
\end{figure}

\subsection{Récapitulatif et demande de devis}
\begin{figure}[H]
  \centering
  \includegraphics[width=\textwidth]{config_recap.png}
  \caption{Étape de finalisation du configurateur}
  \label{fig:config_recap}
\end{figure}

% -----------------------------------------------------------
\section{Authentification}
% -----------------------------------------------------------

\subsection{Formulaire d'inscription}
\begin{figure}[H]
  \centering
  \includegraphics[width=\textwidth]{inscription.png}
  \caption{Authentification — Inscription}
  \label{fig:inscription}
\end{figure}

\subsection{Formulaire de connexion}
\begin{figure}[H]
  \centering
  \includegraphics[width=\textwidth]{connexion.png}
  \caption{Connexion}
  \label{fig:connexion}
\end{figure}

% -----------------------------------------------------------
\section{Espace administrateur}
% -----------------------------------------------------------

\subsection{Tableau de bord administrateur}
\begin{figure}[H]
  \centering
  \includegraphics[width=\textwidth]{admin_dashboard.png}
  \caption{Tableau de bord administrateur}
  \label{fig:admin_dashboard}
\end{figure}

\subsection{Gestion des commandes}
\begin{figure}[H]
  \centering
  \includegraphics[width=\textwidth]{admin_commandes.png}
  \caption{Gestion des commandes}
  \label{fig:admin_commandes}
\end{figure}

\subsection{Gestion des demandes du configurateur}
\begin{figure}[H]
  \centering
  \includegraphics[width=\textwidth]{admin_devis.png}
  \caption{Gestion des demandes du configurateur}
  \label{fig:admin_devis}
\end{figure}

\subsection{Gestion des messages de contact}
\begin{figure}[H]
  \centering
  \includegraphics[width=\textwidth]{admin_messages.png}
  \caption{Gestion des messages de contact}
  \label{fig:admin_messages}
\end{figure}

\subsection{Gestion des utilisateurs}
\begin{figure}[H]
  \centering
  \includegraphics[width=\textwidth]{admin_users.png}
  \caption{Gestion des utilisateurs}
  \label{fig:admin_users}
\end{figure}

\subsection{Gestion des stocks}
\begin{figure}[H]
  \centering
  \includegraphics[width=\textwidth]{admin_stocks.png}
  \caption{Gestion des stocks}
  \label{fig:admin_stocks}
\end{figure}

\subsection{Gestion de sécurité}
\begin{figure}[H]
  \centering
  \includegraphics[width=\textwidth]{admin_securite.png}
  \caption{Gestion de sécurité}
  \label{fig:admin_securite}
\end{figure}

% -----------------------------------------------------------
\section{Espace client}
% -----------------------------------------------------------

\subsection{Section Mon profil}
\begin{figure}[H]
  \centering
  \includegraphics[width=\textwidth]{client_profil.png}
  \caption{Vue de profil du compte client}
  \label{fig:client_profil}
\end{figure}

\subsection{Section "Mes commandes"}
\begin{figure}[H]
  \centering
  \includegraphics[width=\textwidth]{client_commandes.png}
  \caption{Gestion de mes commandes}
  \label{fig:client_commandes}
\end{figure}

\subsection{Section favoris}
\begin{figure}[H]
  \centering
  \includegraphics[width=\textwidth]{client_favoris.png}
  \caption{Gestion des favoris}
  \label{fig:client_favoris}
\end{figure}

\subsection{Section sécurité}
\begin{figure}[H]
  \centering
  \includegraphics[width=\textwidth]{client_securite.png}
  \caption{Gestion de sécurité compte client}
  \label{fig:client_securite}
\end{figure}

% -----------------------------------------------------------
\section{Internationalisation et agent conversationnel}
% -----------------------------------------------------------

\subsection{Internationalisation (FR / AR / EN)}
% [Copier depuis Canva]
\begin{figure}[H]
  \centering
  \includegraphics[width=\textwidth]{i18n.png}
  \caption{Interfaces multilingues}
  \label{fig:i18n}
\end{figure}

\subsection{Agent chatbot intelligent}
% [Copier depuis Canva — mentionner règles conditionnelles, regex, normalisation, RTL]
\begin{figure}[H]
  \centering
  \includegraphics[width=0.5\textwidth]{chatbot.png}
  \caption{Agent chatbot}
  \label{fig:chatbot}
\end{figure}

\subsection{Référencement naturel (SEO)}
% [Copier depuis Canva — meta dynamiques, Open Graph, Schema.org, sitemap.xml]

% -----------------------------------------------------------
\section{Évaluation de l'expérience utilisateur (SUS)}
% -----------------------------------------------------------

La qualité de l'expérience utilisateur de la plateforme FENDRI a été évaluée à l'aide du questionnaire \textbf{SUS (System Usability Scale)}, un outil standardisé composé de 10 questions notées de 1 à 5. Le questionnaire a été soumis à un panel de [N] utilisateurs testeurs.

\begin{table}[H]
\centering
\caption{Résultats de l'évaluation SUS}
\label{tab:sus}
\begin{tabularx}{\textwidth}{c X c}
\toprule
\rowcolor{beige}
\textbf{Participant} & \textbf{Profil} & \textbf{Score SUS /100} \\
\midrule
1  & [Profil] & [Score] \\
2  & [Profil] & [Score] \\
3  & [Profil] & [Score] \\
4  & [Profil] & [Score] \\
5  & [Profil] & [Score] \\
\midrule
\multicolumn{2}{r}{\textbf{Score moyen}} & \textbf{[Moyenne]} \\
\bottomrule
\end{tabularx}
\end{table}

\noindent Un score de \textbf{[X]/100} correspond à une expérience jugée \textbf{[bonne/excellente]} selon l'échelle de référence (Bangor et al., 2008). Les retours qualitatifs ont mis en avant la fluidité du configurateur et la clarté de la navigation.

% -----------------------------------------------------------
\section{Éco-conception numérique}
% -----------------------------------------------------------

La plateforme FENDRI intègre plusieurs bonnes pratiques d'éco-conception numérique visant à réduire son empreinte environnementale tout en améliorant les performances globales.

\subsection*{Réduction de la bande passante — Format WebP}
L'ensemble des visuels produits a été converti au format WebP, permettant une réduction du poids des fichiers de 67 à 70\,\% par rapport aux formats PNG et JPEG traditionnels. Cette optimisation réduit significativement la consommation de données réseau.

\subsection*{Économie d'énergie — Mode sombre}
L'interface du configurateur adopte un thème sombre (fond \texttt{\#0a0a0a}). Sur les écrans OLED et AMOLED, les pixels noirs ne consomment pas d'énergie, ce qui réduit la consommation de la batterie de l'appareil de l'utilisateur.

\subsection*{Architecture légère}
L'architecture repose sur des principes de légèreté : chargement différé des images (\textit{lazy loading}), absence de moteur WebGL, et optimisation des assets via Vite.

\begin{table}[H]
\centering
\caption{Pratiques d'éco-conception intégrées à la plateforme FENDRI}
\label{tab:eco}
\begin{tabularx}{\textwidth}{>{\bfseries}p{4cm} X X}
\toprule
\rowcolor{beige}
\textbf{Pratique} & \textbf{Mesure} & \textbf{Impact} \\
\midrule
Format WebP         & Réduction 67–70\,\% du poids des images & $\downarrow$ Bande passante et consommation réseau \\
Thème sombre        & Fond \texttt{\#0a0a0a}                  & $\downarrow$ Consommation batterie (OLED) \\
Lazy loading        & Chargement à la demande                 & $\downarrow$ Requêtes inutiles au chargement \\
Absence de WebGL    & Images statiques WebP                   & $\downarrow$ Charge GPU, compatibilité universelle \\
Build Vite optimisé & Code-splitting, minification             & $\downarrow$ Taille du bundle JavaScript \\
\bottomrule
\end{tabularx}
\end{table}

\section*{Conclusion}
\addcontentsline{toc}{section}{Conclusion}
Ce chapitre a présenté l'ensemble des interfaces de l'application Domaine FENDRI, en couvrant les parcours utilisateur principaux. L'application répond aux besoins fonctionnels identifiés lors de l'analyse des besoins, en offrant une expérience moderne, multilingue et sécurisée.
